% File: ending.tex
\chapter*{Lời kết}
\addcontentsline{toc}{chapter}{Lời kết} 
\markboth{Lời kết}{} 

\begin{center}
	\vspace*{4cm}
	\rule{0.8\textwidth}{0.4pt} \\[0.8cm]
	\parbox{0.9\textwidth}{
		\centering
		\textit{
			Thế là ta đã đi hết một vòng Thị giác Máy tính. Xin chúc mừng bạn!
		}
	} \\[1cm]
	
	\rule{0.8\textwidth}{0.4pt} \\[2cm]
\end{center}

\noindent
\textbf{Con đường phía trước}

\paragraph{Phần 1}  
Đừng coi “Bách khoa Toàn thư về Lý thuyết” là thứ đọc xong rồi xếp xó.  
Nó chính là nền tảng để bạn dựng mọi mô hình CV sau này.  
Khi có một paper mới ra với tên nghe như phép thuật, quay lại phần này:  
bạn sẽ thấy bản chất nó vẫn xoay quanh convolution, attention, embedding,  
và vài định luật tối ưu quen thuộc. Muốn làm pro thật sự thì không chỉ biết xài model hub,  
mà còn phải hiểu tại sao ResNet lại bền, tại sao Transformer lại bá.  

\paragraph{Phần 2}  
“Cẩm nang Thực chiến” đưa công thức cho bạn, nhưng công thức thì sinh ra để phá.  
Hãy nghịch: thay backbone, thay augment, chỉnh loss, thậm chí làm sai cũng được.  
Mỗi lần model training nổ `CUDA out of memory` hay mAP rớt thảm, hãy coi đó là một mentor thầm lặng.  
Không giáo trình nào dạy nhanh bằng cách tự fix một đống bug và log.

\paragraph{Cuối cùng, đừng quên trách nhiệm.}  
CV không chỉ là code chạy đúng hay accuracy lên đẹp.  
Những mô hình bạn build hôm nay có thể ảnh hưởng tới cách xe tự lái phản ứng,  
cách hệ thống giám sát phân tích camera, hay cách con người tương tác với AI qua hình ảnh.  
Thế nên, ngoài chuyện tối ưu cho nhanh, hãy luôn nghĩ đến bias, fairness và impact.  
Làm kỹ sư CV nên có trách nhiệm đi kèm.

\vfill

\begin{center}
	\parbox{0.9\textwidth}{
		\centering
		{\itshape
			Cảm ơn bạn đã đồng hành tới cuối chặng đường này.  
			Cuốn sách thì đóng lại, nhưng bug thì vẫn còn nhiều, và hành trình học của bạn mới chỉ bắt đầu thôi.  
			Chúc bạn vừa code giỏi, vừa annotate ít lỗi, và vẫn ngủ đủ giấc!
		}
	}
\end{center}
